% Part: many-valued-logic
% Chapter: syntax-and-semantics
% Section: sublogics

\documentclass[../../../include/open-logic-section]{subfiles}

\begin{document}

\olfileid{mvl}{syn}{sub}

\olsection{المنطقات المتعددة القيم بوصفها منطقات جزئية من~$\LogCL$}

المصفوفات التي تعرف بها المنطقات المتعددة القيم المعتادة متفقة على
أمر: متى اقتصرت وسائط دوال الصدق على~$\{\True, \False\}$، جاءت
قيمها على وفق الدوال الكلاسيكية. وتفصيل ذلك أنه إذا كان
$x \in \{\True, \False\}$، كان
$\tf{\lnot}[\Log L](x) = \tf{\lnot}[\LogCL](x)$؛ وإذا كان
$\star$ واحدًا من~$\land$ و$\lor$ و$\lif$، وكان
$x, y \in \{\True, \False\}$، كان
$\tf{\star}[\Log L](x,y) = \tf{\star}[\LogCL](x,y)$. أي إن تقييد
دوال صدق~$\lnot$ و$\land$ و$\lor$ و$\lif$ إلى
$\{\True,\False\}$ يردها بعينها إلى دوال الصدق الكلاسيكية.

\begin{prop}\ollabel{prop:mvl-cl}
  ليكن~$\Log L$ منطقًا متعدد القيم تحتوي لغته~$\lnot$ و$\land$
  و$\lor$ و$\lif$، ولتكن~$\True, \False \in V$، ولتحقق دوال صدقه:
  \begin{enumerate}
    \item\ollabel{prop:not} $\tf{\lnot}[\Log L](x) = \tf{\lnot}[\LogCL](x)$ إذا كان~$x =
    \True$ أو~$x = \False$؛
    \item\ollabel{prop:land} $\tf{\land}[\Log L](x,y) = \tf{\land}[\LogCL](x,y)$،
    \item\ollabel{prop:lor} $\tf{\lor}[\Log L](x,y) = \tf{\lor}[\LogCL](x,y)$،
    \item\ollabel{prop:lif} $\tf{\lif}[\Log L](x,y) = \tf{\lif}[\LogCL](x,y)$،
    إذا كان~$x, y \in \{\True, \False\}$.
  \end{enumerate}
  فعندئذ، لكل تقييم~$\pAssign v$ في~$V$ يحقق
  $\pAssign v(p) \in \{\True,\False\}$ لكل~$p$ واقع في~$!A$، ولكل
  صيغة~$!A$ في الجزء المتولد من متغيرات القضايا والروابط الأربع
  $\lnot$ و$\land$ و$\lor$ و$\lif$، يكون
  $\pValue v[\Log L](!A) = \pValue v[\LogCL](!A)$.
\end{prop}

\begin{proof}
  البرهان بالاستقراء على~$!A$.
  \begin{enumerate}
  \item إذا كانت~$!A \ident p$ ذرية، كان
  $\pValue v[\Log L](!A) = \pAssign v(p) = \pValue
  v[\LogCL](!A)$.
  \item وإذا كان~$!A \ident \lnot B$، حسبنا
  \begin{align*}
    \pValue v[\Log L](!A) & = \tf{\lnot}[\Log L](\pValue v[\Log L](!B)) & &\text{بحسب \olref[val]{defn:pValue}}\\
    & = \tf{\lnot}[\Log L](\pValue v[\LogCL](!B)) && \text{بفرض الاستقراء}\\
    & = \tf{\lnot}[\LogCL](\pValue v[\LogCL](!B))
&&\text{بحسب الفرض \olref{prop:not}،}\\
&&&\text{لأن~$\pValue v[\LogCL](!B) \in \{\True, \False\}$،}\\
& = \pValue v[\LogCL](!A) &&\text{بحسب \olref[val]{defn:pValue}}.
\end{align*}
\item وإذا كان~$!A \ident (!B \land !C)$، كان
\begin{align*}
  \pValue v[\Log L](!A) & = \tf{\land}[\Log L](\pValue v[\Log L](!B), \pValue v[\Log L](!C)) & &\text{بحسب \olref[val]{defn:pValue}}\\
  & = \tf{\land}[\Log L](\pValue v[\LogCL](!B),\pValue v[\LogCL](!C)) && \text{بفرض الاستقراء}\\
  & = \tf{\land}[\LogCL](\pValue v[\LogCL](!B),\pValue v[\LogCL](!C))
&&\text{بحسب الفرض \olref{prop:land}،}\\
&&&\text{لأن~$\pValue v[\LogCL](!B),\pValue v[\LogCL](!C) \in \{\True, \False\}$،}\\
& = \pValue v[\LogCL](!A) &&\text{بحسب \olref[val]{defn:pValue}}.
\end{align*}
\end{enumerate}
وأما حالتا~$!A \ident (!B \lor !C)$ و$!A \ident (!B \lif !C)$
فتجريان على المثال نفسه.
\end{proof}

\begin{cor}
  إذا استوفى منطق متعدد القيم شروط \olref{prop:mvl-cl}، وكان
  $\True \in V^+$ و$\False \notin V^+$، وكانت صيغ
  $\Gamma \cup \{!B\}$ كلها في الجزء المتولد من متغيرات القضايا
  والروابط الأربع المذكورة، لزم من~$\Gamma \Entails[\Log L] !B$
  أن~$\Gamma \Entails[\LogCL] !B$. وعلى الخصوص، كل تحصيل في هذا
  الجزء من~$\Log L$ تحصيل كلاسيكي أيضًا.
\end{cor}

\begin{proof}
  نبرهن المعاكس النقيض. افترض~$\Gamma \Entails/[\LogCL] !B$.
  فهناك~!!{valuation}~$\pAssign v\colon \PVar \to
  \{\True, \False\}$ يحقق~$\pValue v[\LogCL](!A) = \True$ لكل
  $!A \in \Gamma$، ويحقق~$\pValue v[\LogCL](!B) = \False$.
  ولما كان~$\True, \False \in V$، كان~!!{valuation}~$\pAssign v$
  كذلك~!!a{valuation} للمنطق~$\Log L$. ثم تفيد \olref{prop:mvl-cl}
  أن~$\pValue v[\Log L](!A) = \True$ لكل~$!A \in \Gamma$، وأن
  $\pValue v[\Log L](!B) = \False$. ومع~$\True \in V^+$ و
  $\False \notin V^+$، فتتحقق~$\pSat{v}{\Gamma}[\Log L]$
  وتتحقق~$\pSat/{v}{!B}[\Log L]$؛ أي
  $\Gamma \Entails/[\Log L] !B$.
  % تصحيح مصدر (0391-I1، 0391-I2، 0391-N1): قُيدت النتيجة بالجزء الذي تغطيه الفروض، واستُعمل رمز الإشباع من غير نفي مزدوج.
\end{proof}
\end{document}
